\documentclass{article}

\usepackage[separate-uncertainty=true]{siunitx}
\usepackage{graphicx}
\usepackage{amsmath}
\usepackage{bm,upgreek}
\usepackage[colorlinks,colorlinks,linkcolor=blue]{hyperref}
\usepackage[capitalise,noabbrev,nameinlink]{cleveref}
\usepackage{caption}
\usepackage{subcaption}
\usepackage{graphicx}
\usepackage{booktabs}

\title{Photoelectric effect and determination of Planck's constant}
\author{Raymond Langehennig}

\begin{document}
\maketitle

\section{Introduction}
When a metal is subject to electromagnetic radiation, the electrons in the metal may be excited by incident photons to the point that they leave the metal in a phenomenon known as the \emph{photoelectric effect}, which was first explained in the first of Einstein's 1905 \textit{annus mirabilis} papers.

\subsection{Rayleigh and Jeans}
Einstein built off the results of Planck, who in turn had corrected the results of Jeans and Rayleigh.

Rayleigh and Jeans were concerned with modeling the behavior of a \textit{blackbody}, a hypothetical objected that absorbs all incident light, which they could approximate with a \textit{Jeans cube}, a hollow, metal cube with a small hole, such that the hole functioned as a blackbody.



Through the application of statistical mechanics and classical electromagnetism, Jeans and Rayleigh had derived the Rayleigh-Jeans law,
\begin{equation*}
    B_T(\lambda) = \frac{2 c k_\text{B} T}{\lambda^4},
\end{equation*}
by 

\section{Materials and Procedure}

\section{Results}

\section{Discussion}

\section{Conclusions}

\end{document}  